\begin{exercises}
  \recommended \item 
    若一个矩阵具有给定的特征多项式，则其极小多项式有哪些可能？
    \begin{exparts*}
      \partsitem $(x-3)^4$
      \partsitem $(x+1)^3(x-4)$
      \partsitem $(x-2)^2(x-5)^2$
      \partsitem  \( (x+3)^2(x-1)(x-2)^2 \)
    \end{exparts*}
    每种可能的次数是多少？
    \begin{answer} 
      Cayley-Hamilton 定理\nearbytheorem{th:CayHam}指出，
      极小多项式必须含有与特征多项式相同的线性因子，
      尽管次数可以更低，但不能是零次。
      \begin{exparts}
        \partsitem 可能的是
          $m_1(x)=x-3$，$m_2(x)=(x-3)^2$，$m_3(x)=(x-3)^3$
          以及 $m_4(x)=(x-3)^4$。
          第一个是一次多项式，第二个是二次的，
          第三个是三次的，第四个是四次的。
        \partsitem 可能的是 $m_1(x)=(x+1)(x-4)$、
          $m_2(x)=(x+1)^2(x-4)$ 以及 $m_3(x)=(x+1)^3(x-4)$。
          第一个是二次多项式，即次数为二。
          第二个次数为三，第三个次数为四。
        \partsitem 我们有 $m_1(x)=(x-2)(x-5)$、$m_2(x)=(x-2)^2(x-5)$、
          $m_3(x)=(x-2)(x-5)^2$ 以及 $m_4(x)=(x-2)^2(x-5)^2$。
          它们分别是二次、三次、三次与四次多项式。
        \partsitem 可能的是 \( m_1(x)=(x+3)(x-1)(x-2) \)、
          \( m_2(x)=(x+3)^2(x-1)(x-2) \)、
          \( m_3(x)=(x+3)(x-1)(x-2)^2 \)
          以及 \( m_4(x)=(x+3)^2(x-1)(x-2)^2 \)。
          $m_1$ 的次数是三，$m_2$ 的次数是四，
          $m_3$ 的次数是四，$m_4$ 的次数是五。
      \end{exparts}
    \end{answer}
  \item 对于下面这个矩阵
    \begin{equation*}
      \begin{mat}
        0 &1 &0 &1 \\
        1 &0 &1 &0 \\
        0 &1 &0 &1 \\
        1 &0 &1 &0 \\
      \end{mat}
    \end{equation*}
    求其特征多项式与极小多项式。
    \begin{answer}
    特征多项式
    \begin{equation*}
      \deter{T-xI}=
      \begin{vmat}
        -x &1  &0  &1 \\
         1 &-x &1  &0 \\
         0 &1  &-x &1 \\
         1 &0  &1  &-x \\
      \end{vmat}
    \end{equation*}
    是 $x^2(x^2-4)$。

    由 Cayley-Hamilton 定理，特征多项式的所有互不相同的根
    也都是极小多项式的根，所以极小多项式的根为 $0$、$2$ 以及~$-2$。
    因此极小多项式要么是 $x(x^2-4)$，要么是 $x^2(x^2-4)$。
    
    求极小多项式的一种办法是，把矩阵代入这两个表达式试算，
    直到得到零矩阵为止。
    \begin{equation*}
      T(T^2-4I)=
      \begin{mat}
        0 &1 &0 &1 \\
        1 &0 &1 &0 \\
        0 &1 &0 &1 \\
        1 &0 &1 &0 \\
      \end{mat}
      \begin{mat}
        -2 &0  &1  &0 \\
        0  &-2 &0  &2 \\
        2  &0  &-2 &0 \\
        0  &2  &0  &-2 \\
      \end{mat}
      =
      \begin{mat}
        0 &0 &0 &0 \\
        0 &0 &0 &0 \\
        0 &0 &0 &0 \\
        0 &0 &0 &0 \\
      \end{mat}
    \end{equation*}
    所以极小多项式是 $x(x^2-4)$。
    \end{answer}
   \item 
     求此矩阵的极小多项式。
     \begin{equation*}
        \begin{mat}[r]
           0  &1  &0  \\
           0  &0  &1  \\
           1  &0  &0
        \end{mat}
     \end{equation*}
     \begin{answer}
       它的特征多项式有复根。
       \begin{equation*}
          \begin{vmat}
                   -x  &1  &0  \\
                    0  &-x &1  \\
                    1  &0  &-x
          \end{vmat}
          =(1-x)\cdot (x-(-\frac{1}{2}+\frac{\sqrt{3}}{2}i))
                \cdot (x-(-\frac{1}{2}-\frac{\sqrt{3}}{2}i))
       \end{equation*}
       由于根互不相同，特征多项式就等于极小多项式。
     \end{answer}
  \recommended \item 
    求每个矩阵的极小多项式。
    \begin{exparts*}
       \partsitem \( \begin{mat}[r]
                   3  &0  &0  \\
                   1  &3  &0  \\
                   0  &0  &4
                \end{mat} \)
       \partsitem \( \begin{mat}[r]
                   3  &0  &0  \\
                   1  &3  &0  \\
                   0  &0  &3
                \end{mat} \)
       \partsitem \( \begin{mat}[r]
                   3  &0  &0  \\
                   1  &3  &0  \\
                   0  &1  &3
                \end{mat} \)
       \partsitem \( \begin{mat}[r]
                   2  &0  &1  \\
                   0  &6  &2  \\
                   0  &0  &2
                \end{mat} \)
       \partsitem \( \begin{mat}[r]
                   2  &2  &1  \\
                   0  &6  &2  \\
                   0  &0  &2
                \end{mat} \)
       \partsitem \( \begin{mat}[r]
                   -1 &4  &0  &0  &0  \\
                    0 &3  &0  &0  &0  \\
                    0 &-4 &-1 &0  &0  \\
                    3 &-9 &-4 &2  &-1 \\
                    1 &5  &4  &1  &4
                \end{mat} \)
    \end{exparts*}
    \begin{answer}
      每种情形我们都将使用\nearbyexample{ex:MinPolyUsingCH}的方法。
      \begin{exparts}
       \partsitem 因为 $T$ 是三角矩阵，所以 $T-xI$ 也是三角矩阵
         \begin{equation*}
           T-xI=
           \begin{mat}
             3-x  &0    &0   \\
             1    &3-x  &0   \\
             0    &0    &4-x
           \end{mat}
         \end{equation*}
         特征多项式容易算出 $c(x)=\deter{T-xI}=(3-x)^2(4-x)=-1\cdot (x-3)^2(x-4)$。
         极小多项式只有两种可能，
         $m_1(x)=(x-3)(x-4)$ 与 $m_2(x)=(x-3)^2(x-4)$。
         （注意特征多项式带有负号，而极小多项式没有，
         因为它必须首项系数为一。）
         由于 $m_1(T)$ 不是零矩阵
         \begin{equation*}
           (T-3I)(T-4I)
           =
           \begin{mat}[r]
             0  &0  &0  \\
             1  &0  &0  \\
             0  &0  &1
           \end{mat}
           \begin{mat}[r]
             -1  &0  &0  \\
              1  &-1 &0  \\
              0  &0  &0
           \end{mat}
           =
           \begin{mat}[r]
             0  &0  &0  \\
            -1  &0  &0  \\
             0  &0  &0
           \end{mat}
         \end{equation*}
         所以极小多项式是 $m(x)=m_2(x)$。
         \begin{multline*}
           (T-3I)^2(T-4I)
           =(T-3I)\cdot\bigl((T-3I)(T-4I)\bigr)              \\
           =
           \begin{mat}
             0  &0  &0  \\
             1  &0  &0  \\
             0  &0  &1
           \end{mat}
           \begin{mat}[r]
              0  &0  &0  \\
             -1  &0  &0  \\
              0  &0  &0
           \end{mat}
           =
           \begin{mat}[r]
             0  &0  &0  \\
             0  &0  &0  \\
             0  &0  &0
           \end{mat}
         \end{multline*}
       \partsitem 与前一（小）题一样，矩阵是三角矩阵这一事实
        使得特征多项式的计算很容易。
         \begin{equation*}
           c(x)=\deter{T-xI}
               =
               \begin{vmat}
                 3-x  &0   &0   \\
                 1    &3-x &0   \\
                 0    &0   &3-x
               \end{vmat}
               =(3-x)^3=-1\cdot (x-3)^3
         \end{equation*}
         极小多项式有三种可能：
         $m_1(x)=(x-3)$、$m_2(x)=(x-3)^2$ 以及 $m_3(x)=(x-3)^3$。
         我们通过计算 $m_1(T)$
         \begin{equation*}
           T-3I=
           \begin{mat}[r]
             0  &0  &0  \\
             1  &0  &0  \\
             0  &0  &0
           \end{mat}
         \end{equation*}
         与 $m_2(T)$ 来解决这个问题。
         \begin{equation*}
           (T-3I)^2=
           \begin{mat}[r]
             0  &0  &0  \\
             1  &0  &0  \\
             0  &0  &0
           \end{mat}
           \begin{mat}[r]
             0  &0  &0  \\
             1  &0  &0  \\
             0  &0  &0
           \end{mat}
           =          
           \begin{mat}[r]
             0  &0  &0  \\
             0  &0  &0  \\
             0  &0  &0
           \end{mat}
         \end{equation*}
         因为 $m_2(T)$ 是零矩阵，所以 $m_2(x)$ 就是极小多项式。
       \partsitem 同样，该矩阵是三角矩阵。
         \begin{equation*}
           c(x)=\deter{T-xI}
               =
               \begin{vmat}
                 3-x  &0   &0   \\
                 1    &3-x &0   \\
                 0    &1   &3-x
               \end{vmat}
               =(3-x)^3=-1\cdot (x-3)^3
         \end{equation*}
         同样，极小多项式有三种可能：
         $m_1(x)=(x-3)$、$m_2(x)=(x-3)^2$ 以及 $m_3(x)=(x-3)^3$。
         我们计算 $m_1(T)$
         \begin{equation*}
           T-3I=
           \begin{mat}[r]
             0  &0  &0  \\
             1  &0  &0  \\
             0  &1  &0
           \end{mat}
         \end{equation*}
         与 $m_2(T)$
         \begin{equation*}
           (T-3I)^2=
           \begin{mat}[r]
             0  &0  &0  \\
             1  &0  &0  \\
             0  &1  &0
           \end{mat}
           \begin{mat}[r]
             0  &0  &0  \\
             1  &0  &0  \\
             0  &1  &0
           \end{mat}
           =          
           \begin{mat}[r]
             0  &0  &0  \\
             0  &0  &0  \\
             1  &0  &0
           \end{mat}
         \end{equation*}
         与 $m_3(T)$。
         \begin{equation*}
           (T-3I)^3
           =(T-3I)^2(T-3I)
           =
           \begin{mat}[r]
             0  &0  &0  \\
             0  &0  &0  \\
             1  &0  &0
           \end{mat}
           \begin{mat}[r]
             0  &0  &0  \\
             1  &0  &0  \\
             0  &1  &0
           \end{mat}
           =          
           \begin{mat}[r]
             0  &0  &0  \\
             0  &0  &0  \\
             0  &0  &0
           \end{mat}
         \end{equation*}
         因此，极小多项式是 $m(x)=m_3(x)=(x-3)^3$。
       \partsitem 这一情形也是三角矩阵，这里是上三角矩阵。
          \begin{equation*}
            c(x)=\deter{T-xI}=
            \begin{vmat}
              2-x  &0   &1     \\
              0    &6-x &2     \\
              0    &0   &2-x
            \end{vmat}
            =(2-x)^2(6-x)=-(x-2)^2(x-6)
          \end{equation*}
          极小多项式有两种可能，
          $m_1(x)=(x-2)(x-6)$ 与 $m_2(x)=(x-2)^2(x-6)$。
          计算表明极小多项式不是 $m_1(x)$。
          \begin{equation*}
            (T-2I)(T-6I)=
            \begin{mat}[r]
              0  &0  &1  \\
              0  &4  &2  \\
              0  &0  &0  
            \end{mat}
            \begin{mat}[r]
              -4  &0  &1  \\
               0  &0  &2  \\
               0  &0  &-4
            \end{mat}
            =
            \begin{mat}[r]
              0  &0  &-4  \\
              0  &0  &0   \\
              0  &0  &0
            \end{mat}
          \end{equation*}
          因此必有 $m(x)=m_2(x)=(x-2)^2(x-6)$。 
          下面给出验证。
          \begin{multline*}
            (T-2I)^2(T-6I)=(T-2I)\cdot\bigl((T-2I)(T-6I)\bigr)             \\
            =
            \begin{mat}[r]
              0  &0  &1  \\
              0  &4  &2  \\
              0  &0  &0  
            \end{mat}
            \begin{mat}[r]
               0  &0  &-4   \\
               0  &0  &0   \\
               0  &0  &0
            \end{mat}
            =
            \begin{mat}[r]
              0  &0  &0  \\
              0  &0  &0   \\
              0  &0  &0
            \end{mat}
          \end{multline*}
       \partsitem 特征多项式是 
          \begin{equation*}
            c(x)=\deter{T-xI}=
            \begin{vmat}
              2-x  &2   &1     \\
              0    &6-x &2     \\
              0    &0   &2-x
            \end{vmat}
            =(2-x)^2(6-x)=-(x-2)^2(x-6)
          \end{equation*}
          极小多项式有两种可能，
          $m_1(x)=(x-2)(x-6)$ 与 $m_2(x)=(x-2)^2(x-6)$。
          检验第一个
          \begin{equation*}
            (T-2I)(T-6I)=
            \begin{mat}[r]
              0  &2  &1  \\
              0  &4  &2  \\
              0  &0  &0  
            \end{mat}
            \begin{mat}[r]
              -4  &2  &1  \\
               0  &0  &2  \\
               0  &0  &-4
            \end{mat}
            =
            \begin{mat}[r]
              0  &0  &0  \\
              0  &0  &0   \\
              0  &0  &0
            \end{mat}
          \end{equation*}
          这表明极小多项式是
          $m(x)=m_1(x)=(x-2)(x-6)$。
       \partsitem 特征多项式如下。
          \begin{equation*}
            c(x)=\deter{T-xI}=
            \begin{vmat} 
               -1-x &4    &0    &0    &0    \\
                0   &3-x  &0    &0    &0    \\
                0   &-4   &-1-x &0    &0    \\
                3   &-9   &-4   &2-x  &-1   \\
                1   &5    &4    &1    &4-x
            \end{vmat}     
            =(x-3)^3(x+1)^2
          \end{equation*}
          下面是极小多项式的各种可能，按次数升序列出：
          $m_1(x)=(x-3)(x+1)$，$m_1(x)=(x-3)^2(x+1)$，$m_1(x)=(x-3)(x+1)^2$， 
          $m_1(x)=(x-3)^3(x+1)$，$m_1(x)=(x-3)^2(x+1)^2$， 
          以及 $m_1(x)=(x-3)^3(x+1)^2$。 
          第一个行不通
          \begin{align*}
            (T-3I)(T+1I)
            &=
            \begin{mat}[r] 
               -4   &4    &0    &0    &0    \\
                0   &0    &0    &0    &0    \\
                0   &-4   &-4   &0    &0    \\
                3   &-9   &-4   &-1   &-1   \\
                1   &5    &4    &1    &1  
            \end{mat}     
            \begin{mat}[r] 
                0   &4    &0    &0    &0    \\
                0   &4    &0    &0    &0    \\
                0   &-4   &0    &0    &0    \\
                3   &-9   &-4   &3    &-1   \\
                1   &5    &4    &1    &5  
            \end{mat}                           \\     
            &=
            \begin{mat}[r] 
                0   &0    &0    &0    &0    \\
                0   &0    &0    &0    &0    \\
                0   &0    &0    &0    &0    \\
               -4   &-4   &0    &-4   &-4   \\
                4   &4    &0    &4    &4  
            \end{mat}           
          \end{align*}
          而第二个则行得通。
          \begin{multline*}
            (T-3I)^2(T+1I)=(T-3I)\bigl((T-3I)(T+1I)\bigr) \\
            \begin{aligned}
            &=
            \begin{mat}[r] 
               -4   &4    &0    &0    &0    \\
                0   &0    &0    &0    &0    \\
                0   &-4   &-4   &0    &0    \\
                3   &-9   &-4   &-1   &-1   \\
                1   &5    &4    &1    &1  
            \end{mat}     
            \begin{mat}[r] 
                0   &0    &0    &0    &0    \\
                0   &0    &0    &0    &0    \\
                0   &0    &0    &0    &0    \\
               -4   &-4   &0    &-4   &-4   \\
                4   &4    &0    &4    &4  
            \end{mat}                          \\          
            &=
            \begin{mat}[r] 
                0   &0    &0    &0    &0    \\
                0   &0    &0    &0    &0    \\
                0   &0    &0    &0    &0    \\
                0   &0    &0    &0    &0    \\
                0   &0    &0    &0    &0  
            \end{mat}  
            \end{aligned}         
          \end{multline*}
          极小多项式是 \( m(x)=(x-3)^2(x+1) \)。
       \end{exparts} 
     \end{answer}
  \recommended \item 
     在 \( \polyspace_n \) 上，求导算子 $d/dx$ 的极小多项式是什么？
     \begin{answer}
       我们知道 $\polyspace_n$ 是一个维数为 $n+1$ 的空间，
       并且求导算子是幂零指数为~$n+1$ 的幂零算子
       （例如取 $n=3$， 
       $\polyspace_3=\set{c_3x^3+c_2x^2+c_1x+c_0\suchthat c_3,\ldots,c_0\in\C}$，
       而三次多项式的四阶导数是零多项式）。  
       用幂零变换的标准 
       形式来表示这个算子。
       \begin{equation*}
         \begin{mat}
           0  &0  &0  &\ldots &  &0  \\
           1  &0  &0  &       &  &0  \\
           0  &1  &0  &       &  &   \\
              &   &\ddots            \\
           0  &0  &0  &       &1 &0 
         \end{mat}
       \end{equation*}
       这是一个 $\nbyn{(n+1)}$ 矩阵，其特征多项式容易算出，
       为 $c(x)=x^{n+1}$。
       （\textit{注：}这个矩阵是 $\rep{d/dx}{B,B}$，其中
        $B=\sequence{x^n,nx^{n-1},n(n-1)x^{n-2},\ldots,n!}$。）
       为像\nearbyexample{ex:MinPolyUsingCH}那样求出极小多项式，
       我们考虑 $T-0I=T$ 的各次幂。
       但是，当然，$T$ 的第一次变成零矩阵的幂 
       是 $n+1$ 次幂。
       所以极小多项式也是 \( x^{n+1} \)。
     \end{answer}
  \recommended \item 
    求如下形式矩阵的极小多项式
    \begin{equation*}
      \begin{mat}
        \lambda  &0        &0          &\ldots  &        &0  \\
        1        &\lambda  &0          &        &        &0  \\
        0        &1        &\lambda                          \\
                 &         &           &\ddots                \\
                 &         &           &        &\lambda &0   \\
        0        &0        &\ldots     &        &1       &\lambda
      \end{mat}
    \end{equation*}
    其中标量 $\lambda$ 是固定的（即不是一个变量）。
    \begin{answer}
      记该矩阵为 $T$，并设它是 \( \nbyn{n} \) 的。
      由于 $T$ 是三角矩阵，从而 $T-xI$ 也是三角矩阵，
      所以特征多项式是 $c(x)=(x-\lambda)^n$。
      为看出极小多项式与此相同，考虑
      $T-\lambda I$。
      \begin{equation*}
        \begin{mat}
          0        &0        &0          &\ldots  &0  \\
          1        &0        &0          &\ldots  &0  \\
          0        &1        &0                       \\
                   &         &\ddots                  \\
          0        &0        &\ldots     &1       &0      
        \end{mat}
      \end{equation*}
      将它识别为一个幂零指数为~$n$ 的变换的标准形；
      $(T-\lambda I)^j$ 第一次变为零是在 $j$ 等于 $n$ 时。
    \end{answer}
  \item 
    把 \( \polyspace_n \) 中的 \( p(x) \) 映到 \( p(x+1) \) 的
    变换，其极小多项式是什么？
    \begin{answer}
      $n=3$ 的情形提供了一个提示。
      $\polyspace_3$ 的一个自然基是  
      $B=\sequence{1,x,x^2,x^3}$。
      该变换的作用是
      \begin{equation*}
        1\mapsto 1
        \quad
        x\mapsto x+1
        \quad
        x^2\mapsto x^2+2x+1
        \quad
        x^3\mapsto x^3+3x^2+3x+1
      \end{equation*}
      所以表示 $\rep{t}{B,B}$ 是这个上三角矩阵。
      \begin{equation*}
        \begin{mat}[r]
          1  &1  &1  &1  \\
          0  &1  &2  &3  \\
          0  &0  &1  &3  \\
          0  &0  &0  &1
        \end{mat}
      \end{equation*}
      由于它是三角矩阵，显然特征多项式是
      $c(x)=(x-1)^4$。
      至于极小多项式，候选者是 $m_1(x)=(x-1)$、
      \begin{equation*}
        T-1I=
        \begin{mat}[r]
          0  &1  &1  &1  \\
          0  &0  &2  &3  \\
          0  &0  &0  &3  \\
          0  &0  &0  &0
        \end{mat}
      \end{equation*}
      $m_2(x)=(x-1)^2$、 
      \begin{equation*}
        (T-1I)^2=
        \begin{mat}[r]
          0  &0  &2  &6  \\
          0  &0  &0  &6  \\
          0  &0  &0  &0  \\
          0  &0  &0  &0
        \end{mat}
      \end{equation*}
      $m_3(x)=(x-1)^3$、
      \begin{equation*}
        (T-1I)^3=
        \begin{mat}[r]
          0  &0  &0  &6  \\
          0  &0  &0  &0  \\
          0  &0  &0  &0  \\
          0  &0  &0  &0
        \end{mat}
      \end{equation*}
      以及 $m_4(x)=(x-1)^4$。
      由于 $m_1$、$m_2$ 和 $m_3$ 都不对，所以 $m_4$ 必定是对的，
      这也容易验证。
      
      在一般的 $n$ 的情形，表示是一个对角线上全为 $1$ 的
      上三角矩阵。
      于是特征多项式是 $c(x)=(x-1)^{n+1}$。
      验证极小多项式等于特征多项式的一种办法是
      给出类似这样的论证：
      若上三角矩阵的对角线上有非零元素，就称它为 $0$-上三角的；
      若对角线上只有零而对角线上方紧邻处有非零元素，
      就称它为 $1$-上三角的，等等。
      正如上面的例子所说明的，用归纳论证可以证明：
      当 $T$ 只有非负元素时，
      $T^j$ 是 $j$-上三角的。
      % We leave that argument to the reader.
    \end{answer}
   \item 
     把 \( \map{\pi}{\C^3}{\C^3} \) 投影到前两个分量的
     映射，其极小多项式是什么？
      \begin{answer}
        该映射复合两次与复合一次相同：~$\composed{\pi}{\pi}=\pi$，
        即 $\pi^2=\pi$，所以极小多项式的次数
        至多是二，因为 \( m(x)=x^2-x \) 就行得通。
        没有一次多项式行得通这一事实，来自把映射作用到
        $c_1\cdot \pi+c_0\cdot \identity=z$（其中 $z$ 是零映射）
        的左右两边于下面这两个向量。
        \begin{equation*}
          \colvec[r]{0 \\ 0 \\ 1}
          \qquad
          \colvec[r]{1 \\ 0 \\ 0}
        \end{equation*}
        于是极小多项式是 $m$。
      \end{answer}
   \item 
     求一个极小多项式为 \( x^2 \) 的
     \( \nbyn{3} \) 矩阵。
      \begin{answer}
         这是一个答案。
         \begin{equation*}
             \begin{mat}[r]
               0  &0  &0  \\
               1  &0  &0  \\
               0  &0  &0
             \end{mat}
         \end{equation*} 
       \end{answer}
  \item 
     \nearbylemma{le:MatSatItsCharPoly} 的下述声称的证明
     有什么问题：
      ``若 \( c(x)=\deter{T-xI} \) 则 \( c(T)=\deter{T-TI}=0 \,\)''\,？
     \cite{Cullen}
      \begin{answer}
        这里的 \( x \) 必须是标量，而不是矩阵。
      \end{answer}
  \item 
    通过直接计算，对 \( \nbyn{2} \)
    矩阵验证 \nearbylemma{le:MatSatItsCharPoly}。
    \begin{answer}
      矩阵
      \begin{equation*}
         T=\begin{mat}
              a  &b  \\
              c  &d
           \end{mat}
      \end{equation*}
      的特征多项式是 \( (a-x)(d-x)-bc=x^2-(a+d)x+(ad-bc) \)。
      代入
      \begin{multline*}
         \begin{mat}
              a  &b  \\
              c  &d
         \end{mat}^2
         -
         (a+d)\begin{mat}
            a  &b  \\
            c  &d
         \end{mat}
         +
         (ad-bc)\begin{mat}[r]
            1  &0  \\
            0  &1
         \end{mat}                            \\
         =                     
         \begin{mat}
            a^2+bc  &ab+bd  \\
            ac+cd   &bc+d^2
         \end{mat}
         -
         \begin{mat}
            a^2+ad  &ab+bd   \\
            ac+cd   &ad+d^2
         \end{mat}
         +
         \begin{mat}
            ad-bc  &0      \\
            0      &ad-bc
         \end{mat}
      \end{multline*}
      然后只需逐项核对每个元素的和，看结果是不是零矩阵。
    \end{answer}
  \recommended \item
    证明一个 \( \nbyn{n} \) 矩阵的极小多项式的次数
    至多是 \( n \)（而不是人们从小节开头可能猜到的 \( n^2 \)）。
    验证这个最大值 \( n \) 可以取到。
    \begin{answer}
      由 Cayley-Hamilton 定理，极小多项式的次数
      小于或等于特征多项式的次数，
      即 \( n \)。
      \nearbyexample{ex:MinPolyForRotMat} 表明 \( n \) 可以取到。
    \end{answer}
   \recommended \item 
     证明：在非平凡向量空间上，一个线性变换是 
     幂零的当且仅当它仅有的特征值是零。
     \begin{answer}
       设该线性变换为 $\map{t}{V}{V}$。
       若 $t$ 是幂零的，则存在 $n$ 使得 $t^n$ 是零映射，
       于是 $t$ 满足多项式 $p(x)=x^n=(x-0)^n$。
       由 \nearbylemma{le:tSatisImpMinPolyDivides}，$t$ 的极小多项式 
       整除 $p$，所以极小多项式的根只有零。
       由 Cayley-Hamilton 定理，即 \nearbytheorem{th:CayHam}，
       特征多项式的根也只有零。
       于是 $t$ 仅有的特征值是零。

       反之，若 \( n \) 维空间上的变换 \( t \) 只有零这一个特征值， 
       则它的特征多项式是 \( x^n \)。 
       Cayley-Hamilton 定理说映射满足其特征多项式，
       所以 \( t^n \) 是零映射。
       于是 $t$ 是幂零的。
     \end{answer}
   \item 
       零映射或零矩阵的极小多项式是什么？
       恒等映射或单位矩阵的极小多项式又是什么？
       \begin{answer}
         极小多项式必须首项系数为 $1$， 
         因此若一个映射或矩阵的极小多项式 
         是零次多项式，则它将是 $m(x)=1$。
         但恒等映射或单位矩阵等于零映射或零矩阵，
         仅在平凡向量空间上才如此。

         所以在非平凡的情形，极小多项式的次数
         至少是一。
         零映射或零矩阵的极小多项式是 \( m(x)=x \)，而
         恒等映射或单位矩阵的极小多项式是 \( m(x)=x-1 \)。 
       \end{answer}
   \recommended \item 
     对角矩阵的极小多项式是什么？
     \begin{answer}
       对于对角矩阵
       \begin{equation*}
          T=
          \begin{mat}
             t_{1,1}   &0        \\
             0         &t_{2,2}  \\
                       &        &\ddots  \\
                       &        &      &t_{n,n}
          \end{mat}
       \end{equation*}
       特征多项式是 
       $(t_{1,1}-x)(t_{2,2}-x)\cdots (t_{n,n}-x)$。     
       当然，其中一些因子可能重复，例如矩阵可能
       有 $t_{1,1}=t_{2,2}$。
       例如，
       \begin{equation*}
          D=
          \begin{mat}[r]
             3 &0 &0  \\
             0 &3 &0  \\
             0 &0 &1
          \end{mat}
       \end{equation*}
       的特征多项式是 \( (3-x)^2(1-x)=-1\cdot (x-3)^2(x-1) \)。 

       要构造极小多项式， 
       取各项 \( x-t_{i,i} \)，去掉重复项， 
       然后把它们乘起来。
       例如，$D$ 的极小多项式
       是 \( (x-3)(x-1) \)。
       为验证这一点，首先注意 \nearbytheorem{th:CayHam}， 
       即 Cayley-Hamilton 定理，要求特征多项式中的每个线性因子
       在极小多项式中至少出现一次。
       检验另一个方向\Dash 即在对角矩阵的情形， 
       每个线性因子至多出现一次\Dash 的一个办法是
       使用矩阵论证。
       对角矩阵从左边相乘时按
       对角元进行倍乘。
       而在乘积 $(T-t_{1,1}I)\cdots\hbox{}$ 中，即便没有任何重复
       因子，每一行也至少在其中一个因子中为零。 

       例如，在乘积 
       \begin{equation*}
         (D-3I)(D-1I)=(D-3I)(D-1I)I=
         \begin{mat}[r]
           0  &0  &0  \\
           0  &0  &0  \\
           0  &0  &-2        
         \end{mat}
         \begin{mat}[r]
           2  &0  &0  \\
           0  &2  &0  \\
           0  &0  &0
         \end{mat}
         \begin{mat}[r]
           1  &0  &0  \\
           0  &1  &0  \\
           0  &0  &1
         \end{mat}
       \end{equation*}
       中，因为第一个矩阵 $D-3I$ 的第一行与第二行为零，
       所以整个乘积的第一行与第二行都是零。
       又因为中间那个矩阵 $D-1I$ 的第三行为零，
       所以整个乘积的第三行是零。
    \end{answer}
  \recommended \item 
    \definend{投影}\index{projection}%
    \index{transformation!projection} 
    是指满足 \( t^2=t \) 的任一变换 \( t \)。
    （例如，考虑平面 $\Re^2$ 上把每个向量投影到其第一个分量的变换。
    若投影两次，则结果与只投影一次相同。）
    投影的极小多项式是什么？
    \begin{answer}
      本小节开头指出，线性变换的幂不可能永远攀升而没有
      ``重复''，也就是说，对某个幂~$n$ 存在线性关系
      $c_n\cdot t^n+\dots+c_1\cdot t+c_0\cdot \identity=z$，其中 $z$ 是
      零变换。
      投影的定义是，对这样的映射，
      有一个线性关系是二次的：$t^2-t=z$。
      为完成讨论，我们只需考虑这个关系是否可能不是极小的，
      即是否存在极小多项式为常数或一次的投影？

      若极小多项式是常数，则该映射必须满足
      $c_0\cdot\identity=z$，其中 $c_0=1$，因为极小多项式的
      首项系数是 $1$。
      这只在平凡空间上被零变换满足。
      这是一个投影，但不是令人感兴趣的投影。

      若一个变换的极小多项式是一次
      的，则给出 $c_1\cdot t+c_0\cdot\identity=z$，其中 $c_1=1$。
      这个方程给出 $t=-c_0\cdot \identity$。
      再结合要求 $t^2=t$，得
      $t^2=(-c_0)^2\cdot\identity=-c_0\cdot\identity$，由此得
      $c_0=0$ 且 $t$ 是零变换，或者 $c_0=1$ 且
      $t$ 是恒等变换。       

      于是，除投影是零映射或恒等映射的情形外，
      极小多项式是 $m(x)=x^2-x$。 
    \end{answer}
  \item \label{exer:SomeRootsMinPolyAreEigs}
    \textit{本题的前两（小）题是复习。}
    \begin{exparts}
      \partsitem 证明单射映射的复合是单射。
      \partsitem 证明：若线性映射不是单射，则
        定义域中至少有一个非零向量被映到
        陪域中的零向量。
      \partsitem 验证下面摘录的、位于 \nearbytheorem{th:CayHam}
        之前的论断。
         \begin{quotation}
           \noindent \ldots{} 
           若变换 $t$ 的极小多项式 $m(x)$ 
           分解为
           $m(x)=(x-\lambda_1)^{q_1}\cdots (x-\lambda_z)^{q_z}$
           则 
           \( m(t)=\composed{\composed{(t-\lambda_1)^{q_1}}{\cdots}}{
                                                        (t-\lambda_z)^{q_z}} \) 
          是零映射。 
          由于 \( m(t) \) 把每个向量都映为零，所以在映射 \( t-\lambda_i \) 
          中至少有一个把某些
          非零向量映为零。  \ldots{}
          也就是说，\ldots{} 
          至少某些 \( \lambda_i \) 是特征值。
        \end{quotation}
    \end{exparts}
    \begin{answer}
      \begin{exparts}
       \partsitem \textit{这是一般函数的性质，
          而不仅是线性函数的性质。}
          设 $f$ 和 $g$ 是单射函数，且
          $\composed{f}{g}$ 有定义。
          设 $\composed{f}{g}(x_1)=\composed{f}{g}(x_2)$，则
          $f(g(x_1))=f(g(x_2))$。
          由于 $f$ 是单射的，这蕴含 $g(x_1)=g(x_2)$。
          又因 $g$ 也是单射的，这进而蕴含
          $x_1=x_2$。
          于是，总之，$\composed{f}{g}(x_1)=\composed{f}{g}(x_2)$
          蕴含 $x_1=x_2$，所以 $\composed{f}{g}$ 是单射。
        \partsitem 若线性映射 $h$ 
          不是单射，则存在不相等的
          向量 $\vec{v}_1$、$\vec{v}_2$ 被映到相同的值
          $h(\vec{v}_1)=h(\vec{v}_2)$。
          由于 $h$ 是线性的，我们有
          $\zero=h(\vec{v}_1)-h(\vec{v}_2)=h(\vec{v}_1-\vec{v}_2)$，
          于是 $\vec{v}_1-\vec{v}_2$ 是定义域中一个非零向量，
          $h$ 把它映到陪域的零向量  
          （$\vec{v}_1-\vec{v}_2$
          不等于定义域的零向量，因为 $\vec{v}_1$
          不等于 $\vec{v}_2$）。
        \partsitem 极小多项式 
          $m(t)$ 把定义域中的每个向量都映到 
          零，所以它不是单射（在平凡空间中除外，这里 
          不予考虑）。
          由本题的第一（小）题， 
          由于复合 $m(t)$ 不是单射，
          所以分量 $t-\lambda_i$ 中至少有一个不是单射。
          由第二（小）题，$t-\lambda_i$ 有一个非平凡的零空间。
          因为 $(t-\lambda_i)(\vec{v})=\zero$ 成立当且仅当
          $t(\vec{v})=\lambda_i\cdot\vec{v}$，前一句话给出：
          $\lambda_i$ 是特征值（回忆特征值的
          定义要求该关系至少对一个
          非零 $\vec{v}$ 成立）。
      \end{exparts}
    \end{answer}
  \item 
    对或错：~对于 \( n \) 维空间上的变换，若极小多项式的次数为 \( n \)， 
    则该映射可对角化。
    \begin{answer}
      这是错的。
      不可对角化变换的自然例子在这里适用。
      考虑 $\C^2$ 上关于标准基由这个矩阵
      表示的变换。
      \begin{equation*}
        N=
        \begin{mat}[r]
          0  &1  \\
          0  &0
        \end{mat}
      \end{equation*}
      特征多项式是 $c(x)=x^2$。 
      于是极小多项式要么是 $m_1(x)=x$，要么是 $m_2(x)=x^2$。
      第一个不对，因为 $N-0\cdot I$ 不是零矩阵；
      于是在这个例子中极小多项式的次数等于 
      底空间的维数，而且如前所述，
      我们知道这个矩阵不可对角化，因为它是幂零的。 
    \end{answer}
   \item 
     设 $f(x)$ 是一个多项式。
     证明：若 $A$ 与 $B$ 是相似的矩阵，则 $f(A)$ 与 
     $f(B)$ 相似。
     \begin{exparts}
       \partsitem 进而证明相似的矩阵有相同的特征
         多项式。
       \partsitem 证明相似的矩阵有相同的极小多项式。
       \partsitem 判断它们是否相似。
          \begin{equation*}
            \begin{mat}[r]
              1  &3  \\
              2  &3
            \end{mat}
            \qquad
            \begin{mat}[r]
              4  &-1 \\
              1  &1
            \end{mat}
          \end{equation*}
     \end{exparts}
     \begin{answer}
       设 \( A \) 与 \( B \) 相似，即 $A=PBP^{-1}$。
       由这样的事实： 
       \begin{multline*}
           A^n=(PBP^{-1})^n=(PBP^{-1})(PBP^{-1})\cdots(PBP^{-1})   \\
                           =PB(P^{-1}P)B(P^{-1}P)\cdots (P^{-1}P)BP^{-1}
                           =PB^nP^{-1}
       \end{multline*}
       以及 $c\cdot A=c\cdot(PBP^{-1})=P(c\cdot B)P^{-1}$，
       即得所需的事实：对任意多项式函数 $f$，我们有 
       \( f(A)=P\,f(B)\,P^{-1} \)。
       例如，若 $f(x)=x^2+2x+3$，则
       \begin{multline*}
         A^2+2A+3I=(PBP^{-1})^2+2\cdot PBP^{-1}+3\cdot I           \\
                  =(PBP^{-1})(PBP^{-1})+P(2B)P^{-1}+3\cdot PP^{-1}
                  =P(B^2+2B+3I)P^{-1}
       \end{multline*}
       这表明 $f(A)$ 与 $f(B)$ 相似。
       \begin{exparts}
         \partsitem 取 $f$ 为一次多项式，我们得到
           $A-xI$ 与 $B-xI$ 相似。
           相似矩阵有相等的行列式（因为 
           $\deter{A}=\deter{PBP^{-1}}
               =\deter{P}\cdot\deter{B}\cdot\deter{P^{-1}}
               =1\cdot\deter{B}\cdot 1=\deter{B}$）。
           于是特征多项式相等。 
         \partsitem 
           由于 \( P \) 与 \( P^{-1} \) 可逆，\( f(A) \) 是
           零矩阵当且仅当 \( f(B) \) 是零矩阵。
         \partsitem 它们不可能相似，因为它们没有相同的
           特征多项式。
           第一个的特征多项式是 
           $x^2-4x-3$，而第二个的
           特征多项式是 $x^2-5x+5$。
       \end{exparts}  
     \end{answer}
   \item 
    \begin{exparts}
     \partsitem 证明一个矩阵可逆当且仅当 
       其极小多项式中的常数项
       不是 $0$。
     \partsitem 证明：若方阵 \( T \) 不可逆     
       则存在
       非零矩阵 \( S \)，使得 \( ST \) 与 \( TS \) 都等于
       零矩阵。
     \end{exparts}
     \begin{answer}
      设 \( m(x)=x^n+m_{n-1}x^{n-1}+\dots+m_1x+m_0 \)
      是 \( T \) 的极小多项式。
      \begin{exparts}
       \partsitem 
         对于「若」的论证， 
         由于 \( T^n+\dots+m_1T+m_0I \) 是零矩阵，我们
         得到 \( I=(T^n+\dots+m_1T)/(-m_0)=
         T\cdot (T^{n-1}+\dots+m_1I)/(-m_0) \)，所以 
         矩阵 $(-1/m_0)\cdot (T^{n-1}+\dots+m_1I)$ 是 $T$ 的逆。
         对于「仅当」，设 \( m_0=0 \) 
         （我们把 \( n=1 \) 的情形搁置一旁，但它是容易的），
         于是
         \( T^n+\dots+m_1T=(T^{n-1}+\dots+m_1I)T \) 是零矩阵。
         注意 \( T^{n-1}+\dots+m_1I \) 不是零矩阵，因为
         极小多项式的次数是 \( n \)。
         若 \( T^{-1} \) 存在，则把
         \( (T^{n-1}+\dots+m_1I)T \) 与零矩阵从右边
         乘以 $T^{-1}$ 就得出矛盾。
       \partsitem 若 \( T \) 不可逆，则其
         极小多项式中的常数项为零。
         于是，
         \begin{equation*}
            T^n+\dots+m_1T=(T^{n-1}+\dots+m_1I)T=T(T^{n-1}+\dots+m_1I)
         \end{equation*}
         是零矩阵。
       \end{exparts} 
     \end{answer}
   \recommended \item \label{exer:PolyMapsFactor}
      \begin{exparts}
        \partsitem 完成 \nearbylemma{le:PolyMapsFactor} 的证明。
        \partsitem 举一个例子，说明当 $t$ 不是线性时
          该结论不成立。
      \end{exparts}
      \begin{answer}
        \begin{exparts}
          \partsitem
            对于归纳步骤，设 \nearbylemma{le:PolyMapsFactor}
            对次数为 \( i,\ldots,k-1 \) 的多项式
            成立，并考虑次数为 \( k \) 的多项式 \( f(x) \)。 
            把 $f(x)$ 分解为 $f(x)=k(x-\lambda_1)^{q_1}\cdots(x-\lambda_z)^{q_z}$，
            并设
            \( k(x-\lambda_1)^{q_1-1}\cdots(x-\lambda_z)^{q_z} \)
            为 \( c_{n-1}x^{n-1}+\cdots+c_1x+c_0 \)。
            代入：
            \begin{align*}
               k\composed{\composed{(t-\lambda_1)^{q_1}}{\cdots}}{
                                            (t-\lambda_z)^{q_z}}(\vec{v})
               &=
               \composed{(t-\lambda_1)}{
                  \composed{\composed{(t-\lambda_1)^{q_1}}{\cdots}}{
                  (t-\lambda_z)^{q_z}} }
                          (\vec{v})    \\
               &=
               (t-\lambda_1)\,(c_{n-1}t^{n-1}(\vec{v})+\cdots+c_0\vec{v}) \\
               &=
               f(t)(\vec{v})
            \end{align*}
           （第二个等号来自归纳假设， 
           第三个等号来自 \( t \) 的线性）。
         \partsitem 一个例子是考虑平方映射
            $\map{s}{\Re}{\Re}$，由 $s(x)=x^2$ 给出。
            它是非线性的。
            多项式 $f(t)=t^2-1$ 定义的作用 
            把 $s$ 变成 $f(s)=s^2-1$，即这个映射。
            \begin{equation*} 
              x\mapsunder{s^2-1} \composed{s}{s}(x)-1=x^4-1
            \end{equation*}
            注意这个映射不同于映射
            $\composed{(s-1)}{(s+1)}$；例如，第一个映射把
            $x=5$ 变为 $624$，而第二个把 $x=5$ 变为 $675$。
       \end{exparts} 
     \end{answer}
%  \item
%    Give an example of two \( \nbyn{4} \) matrices that have the 
%    same characteristic polynomial and the same minimal polynomial but
%    nonetheless are not similar.
%    \begin{answer}
%      These two have the same characteristic polynomial,
%      $c_S(x)=c_T(x)=(x-2)^4$.
%      \begin{equation*}
%        S=
%        \begin{mat}[r]
%          2  &0  &0  &0  \\
%          1  &2  &0  &0  \\
%          0  &0  &2  &0  \\
%          0  &0  &0  &2
%        \end{mat}
%        \qquad
%        T=
%        \begin{mat}[r]
%          2  &0  &0  &0  \\
%          1  &2  &0  &0  \\
%          0  &0  &2  &0  \\
%          0  &0  &1  &2
%        \end{mat}
%      \end{equation*}
%      For each of the two
%      the candidates for the minimal polynomial are 
%      $m_1(x)=(x-2)$, $m_2(x)=(x-2)^2$, $m_3(x)=(x-2)^3$, and 
%      $m_4(x)=(x-2)^4$.
%      We have 
%      \begin{equation*}
%        S-2I=
%        \begin{mat}[r]
%          0  &0  &0  &0  \\
%          1  &0  &0  &0  \\
%          0  &0  &0  &0  \\
%          0  &0  &0  &0
%        \end{mat}
%        \qquad
%        T-2I=
%        \begin{mat}[r]
%          0  &0  &0  &0  \\
%          1  &0  &0  &0  \\
%          0  &0  &0  &0  \\
%          0  &0  &1  &0
%        \end{mat}
%      \end{equation*}
%      and
%      \begin{equation*}
%        (S-2I)^2=
%        \begin{mat}[r]
%          0  &0  &0  &0  \\
%          0  &0  &0  &0  \\
%          0  &0  &0  &0  \\
%          0  &0  &0  &0
%        \end{mat}
%        =
%        (T-2I)^2
%      \end{equation*}
%      and so the minimal polynomial for each is $m_2$.
%    \end{answer}
  \item
    任何变换或方阵都有极小多项式。
    逆命题成立吗？
    \begin{answer}
      成立。
      按最后一列展开可以检验
      \( x^n+m_{n-1}x^{n-1}+\dots+m_1x+m_0 \) 是下面这个行列式的
      正值或负值。
      \begin{equation*}
         \begin{mat}
            -x  &0  &0  &      &    &m_0    \\
             0  &1-x&0  &      &    &m_1    \\
             0  &0  &1-x&      &    &m_2    \\
                &   &   &\ddots             \\
                &   &   &      &1-x &m_{n-1}
         \end{mat}
      \end{equation*} 
     \end{answer}
\end{exercises}












